\chapter*{\center \Large TRABAJOS FUTUROS}
\addcontentsline{toc}{section}{\bfseries TRABAJOS FUTUROS}
\markboth{TRABAJOS FUTUROS}{TRABAJOS FUTUROS}

A partir de los resultados actuales quedan definidas las siguientes líneas de trabajo:

\begin{itemize}
  \item Evaluación A/B con personas (EXP-040). Pre-registrar y ejecutar el diseño de
        $\sim$1000 juicios (200 por formato) de los cinco formatos frente a las bases U
        (clips equiespaciados con igual presupuesto) y S (mejor señal sola), y decidir si se
        incluye la base T (segmentación de un modelo de lenguaje sobre los subtítulos). Antes
        de medir hay que corregir y declarar el defecto de los destacados en videos de menos
        de 50\,s.
  \item Fusión aprendida como base de comparación. Comparar el gestor con una fusión con
        compuertas aprendidas, como la de A2Summ-X, y no solo con la concatenación, para
        saber si la ventaja se mantiene frente a fusiones más fuertes.
  \item Señales pendientes de extracción. Implementar el texto en pantalla (OCR) y el
        espectro de audio como votantes, cuyos contratos ya están definidos en
        \texttt{signals/}.
  \item Enlaces de trama (G-enlace). Probar descriptores narrativos que capturen continuidad
        de trama y no solo de personajes y situaciones, o reformular la compuerta con jueces
        humanos especializados.
  \item Narrador en documentales. Después de EXP-039, el siguiente candidato es un modelo de
        segmentación con conteo local de hablantes, como pyannote, o una estimación del
        número de hablantes por \emph{eigengap} \cite{park2020nme} sobre ventanas más
        cortas; la anotación propia de los 176 clips de BBC (EXP-035b) sigue pendiente.
  \item Recompensa para políticas de clips. Normalizar la recompensa por posición o
        entrenar con un conjunto grande, dado que la importancia anotada en un conjunto
        pequeño induce preferencias de posición que no transfieren (EXP-038).
  \item AMI y otros dominios. Completar la evaluación de AMI que quedó pendiente por archivos
        de video ausentes y extender el A/B a otros tipos de contenido sin anotación
        pública.
\end{itemize}
